\documentclass[../DoAn.tex]{subfiles}
\begin{document}

\begin{center}
    \Large{\textbf{TÓM TẮT NỘI DUNG ĐỒ ÁN}}\\
\end{center}
\vspace{1cm}
Trong thời đại số hóa hiện nay, việc ứng dụng công nghệ thông tin vào quản lý giáo dục và đào tạo đóng vai trò vô cùng quan trọng trong việc nâng cao hiệu quả công tác quản lý, giảm thiểu thủ tục hành chính và tăng cường tính minh bạch trong quá trình theo dõi, đánh giá. Đặc biệt, đối với các đơn vị quân đội có nhiệm vụ quản lý học viên được gửi đi đào tạo tại các cơ sở giáo dục ngoài Quân đội, việc quản lý thông tin học viên phân tán theo nhiều địa bàn, nhiều trường đào tạo khác nhau đã chứng kiến nhiều thách thức trong việc thu thập, lưu trữ và tổng hợp dữ liệu. Sự phát triển của các hệ thống quản lý tập trung không chỉ mở ra nhiều cơ hội cải thiện hiệu quả quản lý mà còn thay đổi cách thức các đơn vị theo dõi, đánh giá và báo cáo về quá trình đào tạo của học viên.

Mục tiêu của em là nghiên cứu và phát triển một hệ thống phần mềm quản lý học viên gửi đào tạo chuyên biệt, phục vụ công tác quản lý học viên quân đội được Học viện Khoa học Quân sự gửi đi đào tạo tại các trường đại học ngoài Quân đội. Hệ thống tập trung vào việc quản lý toàn diện thông tin học viên, kết quả học tập theo từng học kỳ và năm học, lịch học, thành tích, học phí cùng các thông tin liên quan khác. Đối tượng sử dụng hệ thống bao gồm các cán bộ quản lý, chỉ huy đơn vị và chính bản thân học viên, mỗi nhóm có quyền truy cập và thao tác phù hợp với vai trò của mình. Qua đó, em có thể khám phá các khía cạnh quan trọng của việc xây dựng và vận hành một hệ thống quản lý thông tin hiệu quả, từ việc thiết kế cơ sở dữ liệu, xây dựng giao diện người dùng, quản lý phân quyền và bảo mật, đến việc tích hợp các tính năng tự động hóa như tính toán điểm trung bình tích lũy, tạo lịch cắt cơm tự động dựa trên lịch học. Đây là một bước quan trọng trong việc ứng dụng lý thuyết vào thực tiễn, giúp em có cơ hội trải nghiệm và hiểu rõ hơn về quy trình phát triển một dự án công nghệ thông tin toàn diện, đáp ứng nhu cầu thực tế của công tác quản lý đào tạo. 
\begin{flushright}
Sinh viên thực hiện\\
\begin{tabular}{@{}c@{}}
\textit{(Ký và ghi rõ họ tên)}
\end{tabular}
\end{flushright}

\end{document}